% Compile with: pdflatex report.tex
% Or: latexmk -pdf report.tex
\documentclass[11pt,a4paper]{article}

\usepackage[margin=2.2cm]{geometry}
\usepackage{amsmath}
\usepackage[hidelinks]{hyperref}
\usepackage{graphicx}
\usepackage{booktabs}
\usepackage{tabularx}
\usepackage{xcolor}
\usepackage{enumitem}
\usepackage{parskip}

\setlist[itemize]{topsep=2pt,itemsep=1pt,parsep=0pt}

\title{BrainViz: Dynamic Functional Connectivity Visualization\\\large for the ABIDE Dataset}
\author{}
\date{}

\begin{document}
\maketitle

\section{Introduction}

BrainViz is a web-based tool for exploring dynamic functional connectivity
in resting-state fMRI data from the ABIDE (Autism Brain Imaging Data
Exchange) dataset. It renders time-varying correlation graphs between
14 Resting State Networks (RSNs) as an interactive animation, supporting
both individual subject views and group-level aggregates (average or
median across ASD or HC cohorts). Users can choose between three
connectivity measures (Pearson, Spearman, wavelet coherence) and apply
temporal smoothing or interpolation to shape the output.

\section{Accessing the Application}

The application is live at \url{https://brainviz.ideabattery.org} and
runs in any modern browser (Chrome, Firefox, Safari, Edge). No
installation is required.

\section{Using the Interface}

The UI is organized into collapsible panels on the sidebar, with
playback controls along the bottom.

\paragraph{Data Source.} Choose between two modes:
\begin{itemize}
    \item \textbf{Subject} --- pick a site (from ABIDE I or II), then a
    subject from that site. The diagnosis (ASD or HC) is shown alongside.
    \item \textbf{Collection} --- view an aggregate across all subjects
    of a given group (ASD or HC), using either the \emph{average} or
    \emph{median} across subjects.
\end{itemize}

\paragraph{Correlation.} Select the connectivity measure:
\begin{itemize}
    \item \textbf{Pearson} or \textbf{Spearman} --- sliding-window
    correlation. Configurable window size (5--100) and step (1--100).
    \item \textbf{Wavelet} --- wavelet coherence with phase-lead
    analysis. Offers a toggle between \emph{Both} (show both directions
    of every pair) and \emph{Dominant} (show only the stronger
    direction).
\end{itemize}

\paragraph{Processing.} Optional post-processing of the correlation
time series:
\begin{itemize}
    \item \textbf{Interpolation} --- Linear, Cubic spline, or B-spline,
    at factors 2--10$\times$ to upsample frames.
    \item \textbf{Smoothing} --- Moving average, Exponential, or
    Gaussian, each with a configurable parameter (window, $\alpha$, or
    $\sigma$).
\end{itemize}

\paragraph{Playback.} Playback speed presets (0.5$\times$, 1$\times$,
2$\times$, 4$\times$, 8$\times$); choice of \emph{Linear} or
\emph{Exponential} scale (with adjustable exponent) for mapping edge
weights to color and thickness; and an \textbf{Edge Threshold} slider
that hides edges below the chosen $|w|$.

\paragraph{Nodes.} Per-RSN visibility toggles, with \emph{Hide All} /
\emph{Show All} shortcuts.

\paragraph{Playback controls.} Below the sidebar sit Play/Pause (also
bound to the \texttt{Space} key), \emph{Refresh} (reload data), and
\emph{Download Video} (exports a 4K MP4 of the current animation). A
separate timeline strip along the bottom of the canvas lets the user
scrub through frames.

\section{Data Processing Pipeline}

The pipeline transforms raw per-subject BOLD signals into animated
graph frames in the browser. Two sources feed the backend: per-subject
time-series text files (used for Pearson and Spearman correlation) and
a single pre-computed HDF5 file (used for wavelet coherence). Both are
stored under \texttt{data/} and kept read-only.

\subsection{Raw data sources}

\paragraph{ABIDE I and II time series.} The underlying resting-state
fMRI data comes from the ABIDE I and ABIDE II open datasets. On disk,
subjects are grouped by dataset and acquisition site:
\texttt{data/ABIDE/ABIDE\_\{I,II\}/\textit{site}/dr\_stage1\_\allowbreak
subject\textit{NNNNNNN}.txt}. Each file is the output of a group-ICA
plus dual-regression preprocessing step performed externally: 32
group-level spatial components are estimated across subjects, and
stage-1 dual regression projects them back to yield a
$[T \times 32]$ per-subject time-course matrix (hence the
``dr\_stage1'' prefix). These text files feed the real-time
correlation path (Pearson/Spearman).

\paragraph{Phenotypics.} \texttt{data/phenotypics.csv} maps integer
subject IDs (\texttt{partnum}) to a diagnosis (\texttt{ASD} or
\texttt{HC}). The backend matches filenames such as
\texttt{dr\_stage1\_subject0050649.txt} to CSV rows by parsing the ID
as an integer (stripping leading zeros). Duplicate IDs and missing
diagnoses both raise hard errors at startup.

\paragraph{Wavelet coherence (pre-computed).} Wavelet coherence is too
expensive to compute per request, so it is produced offline in two
stages:

\begin{enumerate}[leftmargin=*]
    \item A MATLAB pipeline (\texttt{data/wavelet\_new/WCMsABIDE.m}
    and \texttt{Wavelet\_coherences\_function.m}) iterates over all 182
    directed RSN pairs. For each pair it loads the per-subject time
    series, computes the Morlet wavelet transform and cross-wavelet
    transform (Grinsted toolbox), derives coherence $R^2$ with Monte
    Carlo significance testing (300 surrogates, 50 scales, $TR = 2$\,s),
    and classifies each $(\text{time}, \text{scale})$ cell of the phase
    angle into one of five categories (Table~\ref{tab:phase}). Output:
    one \texttt{Coherence\_RSN1\_RSN2.mat} file per ordered pair.

    \item A Python conversion script
    (\texttt{backend/scripts/con\-vert\_wavelet\_to\_subject\_files.py})
    reads all \texttt{Coherence\_*.mat} files, validates that their
    subject order matches \texttt{phenotypics.csv} and
    \texttt{participantStructs.mat}, and emits a single
    \texttt{data/wavelet.h5} with the structure below. Only
    \texttt{angle\_maps} (the classified phase codes) is persisted, as
    the leadership ratio computed at runtime is all the frontend needs.
\end{enumerate}

\begin{verbatim}
wavelet.h5
  wavelet_subjects      [n_subjects]            (subject IDs, int)
  period_per_subject    [n_scales, n_subjects]
  pairs/
    aDMN_V1/angle_maps    [n_subjects, n_timepoints, n_scales]
    aDMN_SAL/angle_maps
    ...  (one entry per directed RSN pair)
\end{verbatim}

\subsection{Runtime pipeline}

\begin{enumerate}
    \item \textbf{Raw ABIDE files.} Each subject is a plain-text file
    of shape $[T \times 32]$, where $T$ is the number of fMRI timepoints
    and the columns are the 32 group-ICA components. Repetition time
    (TR) is 2 seconds.

    \item \textbf{RSN filtering.} A fixed subset of 14 of the 32 ICA
    components is retained as the Resting State Networks:
    aDMN, V1, SAL, pDMN, AUD, lFPN, rFPN, latVIS, latSM, CER, SM1, DAN,
    LANG, occVIS.

    \item \textbf{Windowed correlation.} A sliding window of size $W$
    and step $S$ produces a stack of correlation matrices,
    $[F \times 14 \times 14]$ with $F = \lfloor(T-W)/S\rfloor + 1$.
    Pearson/Spearman values lie in $[-1, 1]$; wavelet values lie in
    $[0, 1]$ and represent a leadership ratio (see below).

    \item \textbf{Optional smoothing / interpolation.} Smoothing is
    applied per edge time series; interpolation upsamples the frame
    axis. Both are optional and can be combined.

    \item \textbf{Graph frame construction.} Each correlation matrix
    becomes a graph frame: 14 nodes plus edges. Symmetric methods
    produce 91 edges (upper triangle); the asymmetric wavelet method
    produces up to 182.

    \item \textbf{Rendering.} The frontend lays nodes out on a circle,
    colors and thickens edges according to $|w|$ on a blue $\rightarrow$
    red scale, and filters edges by the user's threshold. The animation
    advances through frames at the selected speed.
\end{enumerate}

\begin{table}[h]
\centering
\small
\begin{tabularx}{\linewidth}{l l l X}
\toprule
Stage & Input shape & Output shape & Key operation \\
\midrule
Parse file          & \texttt{.txt}         & $[T \times 32]$            & \texttt{np.loadtxt} \\
RSN filter          & $[T \times 32]$       & $[T \times 14]$            & Column selection \\
Windowed corr.      & $[T \times 14]$       & $[F \times 14 \times 14]$  & Sliding-window stats \\
Interpolation       & $[F \times 14 \times 14]$ & $[F' \times 14 \times 14]$ & Temporal upsampling \\
Smoothing           & $[F \times 14 \times 14]$ & $[F \times 14 \times 14]$  & Per-edge smoothing \\
Frame build         & $[F \times 14 \times 14]$ & \texttt{GraphFrame[]}  & Matrix $\rightarrow$ nodes + edges \\
Render              & \texttt{GraphFrame}   & Canvas pixels              & d3 scales + canvas API \\
\bottomrule
\end{tabularx}
\caption{Shape transformations through the pipeline.}
\end{table}

\paragraph{Wavelet method (special case).} Wavelet coherence bypasses
steps 1--3: instead of computing correlations on demand, the backend
reads the pre-computed \texttt{wavelet.h5} file described above. At
request time, for each sliding window it counts
\texttt{PHASE\_LEAD} events within an allowed scale band (periods
10--100\,s) and reports a leadership ratio. Both matrix entries
$(i, j)$ and $(j, i)$ are populated, making the graph asymmetric. The
five phase codes produced by the MATLAB pipeline are:

\begin{table}[h]
\centering
\small
\begin{tabularx}{\linewidth}{l l X}
\toprule
Code & Label & Meaning \\
\midrule
$+1$ & \texttt{PHASE\_LEAD}     & Network A leads network B (${\sim}{+}90^\circ$) \\
$-1$ & \texttt{PHASE\_LAG}      & Network A lags network B (${\sim}{-}90^\circ$) \\
$+2$ & \texttt{PHASE\_IN\_PHASE}& Networks are in phase \\
$-2$ & \texttt{PHASE\_ANTI}     & Networks are anti-phase \\
$\phantom{+}0$ & \texttt{PHASE\_NONE}   & No significant coherence (outside COI or $R^2 < 0.5$) \\
\bottomrule
\end{tabularx}
\caption{Wavelet phase-angle categories.}
\label{tab:phase}
\end{table}

The leadership ratio is $n_{\text{lead}} / n_{\text{all}} \in [0, 1]$
over all $(\text{time}, \text{scale})$ cells in the window, after
masking cells where either direction reports \texttt{PHASE\_NONE}.

\end{document}
